\clearpage\thispagestyle{empty}
\pdfbookmark[0]{Sources for Pallas (29 May 1814)}{pallas1814may-sources}
\begin{center}\textbf{Sources for Pallas (29 May 1814)}\end{center}
\noindent This AI translation follows Carl Friedrich Gauss’s letter headed Pallas in his \href{https://gdz.sub.uni-goettingen.de/id/PPN236007467}{\textit{Werke}, vol.~VI} (Göttingen, 1874), pp.~376--377. Gauss sent it from Göttingen on 29 May 1814 to Johann Elert Bode, editor of the \textit{Astronomisches Jahrbuch}. It appeared in the yearbook for 1817 (Berlin, 1814), pp.~\href{https://api.digitale-sammlungen.de/iiif/image/v2/bsb10538346_00220/full/1600,/0/default.jpg}{212}, \href{https://api.digitale-sammlungen.de/iiif/image/v2/bsb10538346_00221/full/1600,/0/default.jpg}{213}, \href{https://api.digitale-sammlungen.de/iiif/image/v2/bsb10538346_00222/full/1600,/0/default.jpg}{214}, \href{https://api.digitale-sammlungen.de/iiif/image/v2/bsb10538346_00223/full/1600,/0/default.jpg}{215}. The links lead to the Bavarian State Library’s digitized copy.
\par\smallskip\noindent The first printing’s heading names Juno’s observations and new orbital elements, the next opposition of Pallas and astronomical news. A note on p.~212 signed B. identifies the earlier Juno letter as dated 10 December 1813. The collected edition prints the letter under Pallas and gives neither the longer heading nor that note. Both editions give the instrument maker’s name as Frauenhofer. Translation and typesetting produced by AI.
\clearpage
